% ==============================================================================
% CHAPITRE 7 : FONCTIONS LINÉAIRES, AFFINES & LECTURES GRAPHIQUES (VERSION ÉLÈVE)
% ==============================================================================
\chapter{Fonctions Linéaires, Affines \& Lectures Graphiques}

% ------------------------------------------------------------------------------
% 1. ACTIVITÉ D'INVESTIGATION (SITUATION PROFESSIONNELLE : TRANSPORT / FORFAIT)
% ------------------------------------------------------------------------------
\begin{activite}{Tarification d'un service de livraison express (Transport \& Logistique)}
Une entreprise de transport propose deux formules tarifaires pour la livraison régionale de colis :
\begin{itemize}[leftmargin=1.5em, itemsep=1pt]
    \item \textbf{Tarif Éco (Formule Linéaire)} : Aucun abonnement, facturation de \textbf{2 \euro{}} par kilomètre parcouru $x$.
    \item \textbf{Tarif Pro (Formule Affine)} : Forfait fixe de prise en charge de \textbf{30 \euro{}}, puis \textbf{1 \euro{}} par kilomètre parcouru $x$.
\end{itemize}

\vspace{1mm}
\noindent
\begin{minipage}[c]{0.48\linewidth}
\centering
\begin{tikzpicture}[scale=0.55]
    % Repère
    \draw[step=1cm, slate200, very thin] (0,0) grid (8,6);
    \draw[thick, ->, >=stealth, slate800] (0,0) -- (8.4,0) node[right, font=\tiny] {$x$ (km)};
    \draw[thick, ->, >=stealth, slate800] (0,0) -- (0,6.4) node[above, font=\tiny] {$y$ (\euro{})};
    
    % Graduations
    \foreach \x/\val in {2/20, 4/40, 6/60, 8/80}
        \draw (\x,0.1) -- (\x,-0.1) node[below, font=\tiny] {\val};
    \foreach \y/\val in {1/20, 2/40, 3/60, 4/80, 5/100, 6/120}
        \draw (0.1,\y) -- (-0.1,\y) node[left, font=\tiny] {\val};
        
    % Droites
    \draw[line width=1.3pt, colPrimary] (0,0) -- (4,4) -- (6,6) node[above left, font=\bfseries\tiny] {Éco : $f(x)=2x$};
    \draw[line width=1.3pt, colSecondary] (0,1.5) -- (6,4.5) -- (8,5.5) node[above right, font=\bfseries\tiny] {Pro : $g(x)=x+30$};
    
    % Point d'intersection
    \fill[amber!90!black] (3,3) circle (3pt);
    \draw[dashed, thick, amber!90!black] (3,0) -- (3,3) -- (0,3);
    \node[above left, font=\bfseries\tiny, amber!90!black] at (3.2,3) {$(30\text{ km} ; 60\text{ \euro{}})$};
\end{tikzpicture}
\end{minipage}\hfill
\begin{minipage}[c]{0.50\linewidth}
\begin{tcolorbox}[colback=slate100, colframe=colPrimary, arc=4pt, boxrule=0.8pt, top=4pt, bottom=4pt, left=6pt, right=6pt]
\sffamily\footnotesize
\textbf{Problématique :}
\begin{enumerate}[leftmargin=1.2em, itemsep=2pt]
    \item Quelle droite passe par l'origine $(0;0)$ ? Que traduit-elle sur la proportionnalité ?
    \item À partir de quelle distance le Tarif Pro devient-il plus économique que le Tarif Éco ?
    \item Exprimer le coût $f(x)$ et $g(x)$ en fonction du nombre de kilomètres $x$.
\end{enumerate}
\end{tcolorbox}
\end{minipage}

\vspace{2mm}
\noindent\textbf{Espace de recherche :}
\lignesreponse[3]
\end{activite}

% ------------------------------------------------------------------------------
% 2. COURS : NOTION GÉNÉRALE DE FONCTION, IMAGE & ANTÉCÉDENT
% ------------------------------------------------------------------------------
\section{Notion de Fonction, Image \& Antécédent}

\begin{cadredef}[Fonction, Image et Antécédent]
Une \textbf{fonction} $f$ est un processus mathématique qui, à chaque nombre de départ noté $x$, associe un \textbf{unique} nombre d'arrivée noté $f(x)$ (se lit « $f$ de $x$ »).
\[ x \xrightarrow{\quad f \quad} f(x) \]
\begin{itemize}[leftmargin=1.5em, itemsep=2pt]
    \item $f(x)$ est l'\textbf{image} de $x$ par la fonction $f$ (lue sur l'axe vertical des ordonnées $y$).
    \item $x$ est un \textbf{antécédent} de $f(x)$ par la fonction $f$ (lu sur l'axe horizontal des abscisses $x$).
\end{itemize}
\end{cadredef}

\begin{cadremethode}[Calculer une Image et un Antécédent]
Soit la fonction $f$ définie par $f(x) = 3x - 5$.
\begin{enumerate}[leftmargin=1.5em, itemsep=2pt]
    \item \textbf{Calculer l'image de 4} : On remplace $x$ par $4$ dans la formule :
    \[ f(4) = 3 \times 4 - 5 = 12 - 5 = \mathbf{7} \qquad \text{(L'image de $4$ est $7$)}. \]
    \item \textbf{Calculer l'antécédent de 16} : On résout l'équation $f(x) = 16$ :
    \[ 3x - 5 = 16 \iff 3x = 16 + 5 \iff 3x = 21 \iff x = \frac{21}{3} = \mathbf{7} \qquad \text{(L'antécédent de $16$ est $7$)}. \]
\end{enumerate}
\end{cadremethode}

\begin{cadreexemple}
Soit la fonction $g(x) = 2x + 8$.
\begin{itemize}[leftmargin=1.5em, itemsep=2pt]
    \item Calculer l'image de $5$ : $g(5) = 2 \times \pointille[0.8cm] + 8 = \pointille[1cm] + 8 = \pointille[1.5cm]$
    \item Trouver l'antécédent de $20$ : $2x + 8 = 20 \iff 2x = 20 - 8 \iff 2x = \pointille[1cm] \iff x = \pointille[1.5cm]$
\end{itemize}
\end{cadreexemple}

% ------------------------------------------------------------------------------
% 3. COURS : FONCTIONS LINÉAIRES & AFFINES
% ------------------------------------------------------------------------------
\section{Fonctions Linéaires \& Fonctions Affines}

\begin{cadredef}[Fonction Linéaire]
Une \textbf{fonction linéaire} est une fonction de la forme :
\[ \mathbf{f(x) = ax} \quad (\text{où } a \text{ est un nombre réel fixé}) \]
\begin{itemize}[leftmargin=1.5em, itemsep=2pt]
    \item Elle traduit une situation de \textbf{proportionnalité}.
    \item Sa représentation graphique est une \textbf{droite passant par l'origine} $(0;0)$.
    \item Le nombre $a$ est appelé le \textbf{coefficient directeur} (la pente de la droite).
\end{itemize}
\end{cadredef}

\begin{cadredef}[Fonction Affine]
Une \textbf{fonction affine} est une fonction de la forme :
\[ \mathbf{f(x) = ax + b} \quad (\text{où } a \text{ et } b \text{ sont des nombres fixés}) \]
\begin{itemize}[leftmargin=1.5em, itemsep=2pt]
    \item Sa représentation graphique est une \textbf{droite}.
    \item $\mathbf{a}$ est le \textbf{coefficient directeur} (inclinaison / pente de la droite).
    \item $\mathbf{b}$ est l'\textbf{ordonnée à l'origine} (point d'intersection avec l'axe vertical $(y)$ en $(0;b)$).
\end{itemize}
\end{cadredef}

\begin{cadremethode}[Tracer la droite d'une fonction affine]
Pour tracer la droite représentant la fonction $f(x) = 2x - 3$ :
\begin{enumerate}[leftmargin=1.5em, itemsep=2pt]
    \item On choisit deux valeurs de $x$ (ex : $x = 0$ et $x = 3$).
    \item On calcule leurs images : $f(0) = -3 \implies A(0\,;\,-3)$ et $f(3) = 3 \implies B(3\,;\,3)$.
    \item On place les points $A$ et $B$ dans le repère, puis on trace la droite à la règle.
\end{enumerate}
\end{cadremethode}

\begin{cadreretenu}[L'Essentiel du Chapitre 7]
\begin{itemize}[leftmargin=1.5em, itemsep=2pt]
    \item \textbf{Image} : calculée sur l'axe vertical ($y$) en remplaçant $x$ dans $f(x)$.
    \item \textbf{Antécédent} : calculé sur l'axe horizontal ($x$) en résolvant l'équation $f(x) = k$.
    \item \textbf{Fonction linéaire} : $f(x) = ax \implies$ droite qui \textbf{passe par l'origine} $(0;0)$.
    \item \textbf{Fonction affine} : $f(x) = ax + b \implies$ droite qui coupe l'axe vertical en $(0;b)$.
    \item Si $a > 0$, la droite monte (croissante). Si $a < 0$, la droite descend (décroissante).
\end{itemize}
\end{cadreretenu}

% ------------------------------------------------------------------------------
% 4. QUESTIONS FLASH / AUTOMATISMES
% ------------------------------------------------------------------------------
\section{Questions Flash / Automatismes}

\begin{automatisme}[8 Questions Flash d'Entraînement]
\begin{enumerate}[label=\textbf{Q\arabic*.}]
    \item Soit $f(x) = 4x$. Calculer l'image de $3$ : $f(3) = \pointille[3cm]$
    \item Soit $g(x) = 5x - 2$. Calculer l'image de $4$ : $g(4) = \pointille[3cm]$
    \item La fonction $h(x) = -3x$ est-elle linéaire ou affine ? \pointille[3cm]
    \item La fonction $k(x) = 2x + 7$ est-elle linéaire ou affine ? \pointille[3cm]
    \item Donner l'ordonnée à l'origine de la droite $y = 3x - 8$ : $b = \pointille[3cm]$
    \item Donner le coefficient directeur de la droite $y = -5x + 4$ : $a = \pointille[3cm]$
    \item Trouver l'antécédent de $15$ par la fonction linéaire $f(x) = 3x$ : $x = \pointille[3cm]$
    \item La droite représentant $f(x) = 7x$ passe-t-elle par l'origine $(0;0)$ ? \pointille[3cm]
\end{enumerate}
\end{automatisme}

% ------------------------------------------------------------------------------
% 5. TRAVAUX DIRIGÉS (TD) - EXERCICES D'ENTRAÎNEMENT
% ------------------------------------------------------------------------------
\section{Travaux Dirigés (TD) : Exercices d'Entraînement}

\begin{exercice}[2 pts]{\capRealiser}{Calculs d'images et tableau de valeurs}
Soit la fonction affine $f$ définie par : $f(x) = 3x - 4$.
\begin{enumerate}[leftmargin=1.5em]
    \item Calculer les images suivantes :
    \begin{itemize}
        \item $f(-2) = 3 \times (-2) - 4 = \pointille[3cm]$
        \item $f(0) = 3 \times 0 - 4 = \pointille[3cm]$
        \item $f(5) = 3 \times 5 - 4 = \pointille[3cm]$
    \end{itemize}
    \item Déterminer par le calcul l'antécédent de $17$ par la fonction $f$.
\end{enumerate}
\lignesreponse[2]
\end{exercice}

\begin{exercice}[3 pts]{\capSApproprier}{Lectures graphiques directes}
On donne la représentation graphique de la fonction $f$ ci-dessous :

\begin{center}
\begin{tikzpicture}[scale=0.7]
    \draw[step=1cm, slate200, very thin] (-3,-2) grid (5,5);
    \draw[thick, ->, >=stealth, slate800] (-3,0) -- (5.3,0) node[right, font=\scriptsize] {$x$};
    \draw[thick, ->, >=stealth, slate800] (0,-2) -- (0,5.3) node[above, font=\scriptsize] {$y$};
    \foreach \x in {-2,-1,1,2,3,4}
        \draw (\x,0.1) -- (\x,-0.1) node[below, font=\tiny] {\x};
    \foreach \y in {-1,1,2,3,4}
        \draw (0.1,\y) -- (-0.1,\y) node[left, font=\tiny] {\y};
    \node[below left, font=\tiny] at (0,0) {$0$};
    
    % Droite y = x + 1
    \draw[line width=1.3pt, colPrimary] (-3,-2) -- (4,5) node[above right, font=\footnotesize] {$(d_f)$};
\end{tikzpicture}
\end{center}

Par lecture graphique :
\begin{enumerate}[leftmargin=1.5em]
    \item Lire l'image de $2$ par $f$ : $f(2) = \pointille[2cm]$
    \item Lire l'image de $-1$ par $f$ : $f(-1) = \pointille[2cm]$
    \item Lire l'antécédent de $4$ par $f$ : $x = \pointille[2cm]$
    \item Donner l'ordonnée à l'origine de cette droite : $b = \pointille[2cm]$
\end{enumerate}
\end{exercice}

\begin{exercice}[3 pts]{\capRealiser}{Tracé de représentations graphiques}
Dans le repère ci-dessous :
\begin{enumerate}[leftmargin=1.5em]
    \item Tracer la droite $(d_1)$ représentant la fonction linéaire $f(x) = 2x$.
    \item Tracer la droite $(d_2)$ représentant la fonction affine $g(x) = -x + 3$.
\end{enumerate}

\begin{center}
\begin{tikzpicture}[scale=0.7]
    \draw[step=1cm, slate200, very thin] (-3,-3) grid (5,5);
    \draw[thick, ->, >=stealth, slate800] (-3,0) -- (5.3,0) node[right, font=\scriptsize] {$x$};
    \draw[thick, ->, >=stealth, slate800] (0,-3) -- (0,5.3) node[above, font=\scriptsize] {$y$};
    \foreach \x in {-2,-1,1,2,3,4}
        \draw (\x,0.1) -- (\x,-0.1) node[below, font=\tiny] {\x};
    \foreach \y in {-2,-1,1,2,3,4}
        \draw (0.1,\y) -- (-0.1,\y) node[left, font=\tiny] {\y};
    \node[below left, font=\tiny] at (0,0) {$0$};
\end{tikzpicture}
\end{center}
\end{exercice}

\begin{exercice}[4 pts]{\capAnalyser}{Comparaison de Tarifs Télécom (Vente / Tertiaire)}
Un opérateur propose deux formules pour les communications professionnelles :
\begin{itemize}[leftmargin=1.5em, itemsep=1pt]
    \item \textbf{Formule A} : $0{,}20\text{ \euro{}}$ par minute de communication.
    \item \textbf{Formule B} : Abonnement fixe de $12\text{ \euro{}}$ par mois plus $0{,}05\text{ \euro{}}$ par minute.
\end{itemize}

\begin{enumerate}[leftmargin=1.5em]
    \item Exprimer le coût mensuel $C_A(x)$ de la Formule A pour $x$ minutes : $C_A(x) = \pointille[3cm]$
    \item Exprimer le coût mensuel $C_B(x)$ de la Formule B pour $x$ minutes : $C_B(x) = \pointille[3cm]$
    \item Préciser la nature de chaque fonction ($C_A$ et $C_B$).
    \item Calculer le coût de chaque formule pour $x = 100\text{ minutes}$ d'appels dans le mois.
    \item Résoudre l'équation $0{,}20x = 0{,}05x + 12$. Interpréter la valeur trouvée.
\end{enumerate}
\lignesreponse[4]
\end{exercice}

\begin{exercice}[4 pts]{\capAnalyser}{Facturation d'Intervention d'un Chauffagiste (BTP / Énergie)}
Un artisan chauffagiste facture ses dépannages à domicile selon le barème suivant :
Forfait déplacement fixe de \textbf{45 \euro{}}, plus un tarif horaire de \textbf{35 \euro{}} par heure passée.

\begin{enumerate}[leftmargin=1.5em]
    \item Exprimer le montant total de la facture $F(x)$ en fonction du nombre d'heures d'intervention $x$.
    \item S'agit-il d'une fonction linéaire ou affine ? Justifier.
    \item Une intervention a duré $3\text{ heures et demie}$ ($3{,}5\text{ h}$). Calculer le montant de la facture.
    \item Un client a payé une facture totale de $185\text{ \euro{}}$. Combien de temps le chauffagiste a-t-il travaillé chez ce client ?
\end{enumerate}
\lignesreponse[4]
\end{exercice}

% ------------------------------------------------------------------------------
% 6. TP TICE / CALCULATRICE NUMWORKS
% ------------------------------------------------------------------------------
\section{TP TICE : Étude de Fonctions sur NumWorks}

\begin{tpinfo}[Calculatrice NumWorks]{Application Fonctions \& Tableau de Valeurs}
\begin{enumerate}[leftmargin=1.5em, itemsep=2pt]
    \item \textbf{Ajouter une fonction (Application Fonctions)} :
    \begin{itemize}
        \item Ouvrir l'application \textbf{Fonctions} (icône bleue avec courbe).
        \item Sélectionner \textbf{Ajouter une fonction} $\rightarrow$ choisir \textbf{Affine} ou taper : \texttt{f(x) = 2x - 3}.
        \item Ajouter une seconde fonction : \texttt{g(x) = -x + 6}.
    \end{itemize}
    \item \textbf{Tracer le graphique} :
    \begin{itemize}
        \item Aller dans l'onglet \textbf{Graphique} en haut de l'écran avec les flèches.
        \item Observer les deux droites et déplacer le curseur au point d'intersection.
        \item Noter les coordonnées du point d'intersection : $(\pointille[1.5cm]\,;\,\pointille[1.5cm])$.
    \end{itemize}
    \item \textbf{Afficher le Tableau de valeurs} :
    \begin{itemize}
        \item Aller dans l'onglet \textbf{Tableau}.
        \item Régler l'intervalle avec \textbf{Régler le tableau} (Début : $0$, Fin : $10$, Pas : $1$).
        \item Noter la valeur de $f(5)$ affichée par la calculatrice : $f(5) = \pointille[2cm]$.
    \end{itemize}
\end{enumerate}
\end{tpinfo}

% ------------------------------------------------------------------------------
% 7. VERS LE BREVET (DNB PRO)
% ------------------------------------------------------------------------------
\section{Vers le Brevet (DNB Pro)}

\begin{sujetdnb}[DNB Pro Métropole][10 pts]{Tarifs d'Entrée à un Parc Aquatique}
Un parc d'attractions aquatique propose deux options d'accès pour la saison estivale :
\begin{itemize}[leftmargin=1.5em, itemsep=1pt]
    \item \textbf{Option 1 (Tarif Plein)} : $15\text{ \euro{}}$ par entrée.
    \item \textbf{Option 2 (Carte Club)} : Achat d'une carte d'adhérent à $40\text{ \euro{}}$ pour la saison, puis $7\text{ \euro{}}$ par entrée.
\end{itemize}

\begin{center}
\renewcommand{\arraystretch}{1.2}
\begin{tabularx}{\linewidth}{|l|X|X|X|}
\hline
\rowcolor{slate800} \textcolor{white}{\textbf{Nombre d'entrées ($x$)}} & \textcolor{white}{\textbf{2 entrées}} & \textcolor{white}{\textbf{5 entrées}} & \textcolor{white}{\textbf{10 entrées}} \\
\hline
\textbf{Prix Option 1 ($P_1(x)$)} & $30\text{ \euro{}}$ & \pointille[2cm] & \pointille[2cm] \\
\hline
\textbf{Prix Option 2 ($P_2(x)$)} & $40 + 2 \times 7 = 54\text{ \euro{}}$ & \pointille[2cm] & \pointille[2cm] \\
\hline
\end{tabularx}
\end{center}

\begin{enumerate}[leftmargin=1.5em, itemsep=2pt]
    \item \capSApproprier\ Compléter le tableau de valeurs ci-dessus.
    \item \capRealiser\ Exprimer $P_1(x)$ et $P_2(x)$ en fonction du nombre d'entrées $x$.
    \item \capAnalyser\ Parmi ces deux fonctions, laquelle est une fonction linéaire ? Justifier.
    \item \capRealiser\ Résoudre l'équation : $15x = 7x + 40$.
    \item \capValider\ Que représente la solution de cette équation pour un visiteur ?
    \item \capCommuniquer\ Un jeune prévoit de se rendre $8\text{ fois}$ au parc aquatique durant l'été. Quelle option doit-il choisir pour payer le moins cher ? Justifier avec les calculs détaillés.
\end{enumerate}

\vspace{1mm}
\noindent\textbf{Rédigez votre démarche ci-dessous :}
\lignesreponse[6]
\end{sujetdnb}
